\subsection{A fixed finite word for the degree-\(163\) obstruction}
\label{sec:163-finite-word}

Use \(R=R_{\rm global}\), \(P=P_{163}\), \(A=Q_{163}P\),
and \(E=P/A\). Put
\[
 \mathcal O=\mathbb Z_2[\zeta_{163}],\quad
 s(\zeta_{163})=\zeta_{163}^2,\quad c=s^{81},\quad
 W=\ker\bigl(\operatorname{Tr}_{\mathbb F_8/\mathbb F_2}\bigr).
\]
The exact detector proved above identifies \(E\) with the
two-dimensional additive space \(W\):
\begin{equation}
 \Psi(t)=\sum_{j=0}^{26}s^{3j}
             \left((t-1)/2\bmod2\mathcal O\right),\qquad
 \ker(\Psi:P\to W)=A.
 \label{eq:163-trace-detector}
\end{equation}
Here \(W\) is a vector space inside \(\mathbb F_8\), not a
subfield of order four.

Let \(\mathcal O_4=\mathcal O/4\mathcal O\), and define the fixed
finite group
\[
 \mathscr D=(\mathcal O_4)_{\rm add}\rtimes\langle t\rangle,
 \qquad t^{163}=1,\qquad tvt^{-1}=\zeta_{163}v.
\]
Write \(d=\zeta_{163}-1\), a unit, and put
\[
 X=(0,t),\qquad Y=(d^{-1},t),\qquad n=2^{81}+1.
\]
For \(0\leq j\leq26\), let \(a_j\) be the least positive residue
of \(2^{3j}\) modulo \(163\).

\begin{proposition}
\label{prop:163-finite-word-detector}
The elements \(X,Y\) generate \(\mathscr D\), so they define
one finite quotient of \(\widehat F_2\), of order \(163\cdot4^{162}\).
For every \(r\in R\), use the conjugator \(f_{r^n}\) of the
group power \(r^n\), in the convention
\[
 \phi_f(x)=x,\qquad \phi_f(y)=f^{-1}yf,\qquad
 [x,y]=xyx^{-1}y^{-1}.
\]
Then
\begin{equation}
 \prod_{j=0}^{26}f_{r^n}(X^{a_j},Y^{a_j})
       =\bigl(2d^{-2}\Psi(h_{163}(r)),1\bigr)
       \quad\text{in }\mathscr D.
 \label{eq:163-exact-word-value}
\end{equation}
In particular,
\begin{equation}
 \prod_{j=0}^{26}f_{r^n}(X^{a_j},Y^{a_j})=1
       \quad\Longleftrightarrow\quad h_{163}(r)\in A.
 \label{eq:163-word-equivalence}
\end{equation}
The word value lies in \(2d^{-2}W\), a subgroup of order four.
\end{proposition}
\begin{proof}
The multiplication is
\((u,t^i)(v,t^j)=(u+\zeta_{163}^iv,t^{i+j})\).
It gives \([X,Y]=(1,1)\).
Conjugating this translation by powers of \(X\) gives
\((\zeta_{163}^i,1)\). These elements span \(\mathcal O_4\)
over \(\mathbb Z/4\mathbb Z\), and \(X\) generates the quotient
of order \(163\). Hence \(X,Y\) generate \(\mathscr D\).
Its derived subgroup is exactly the translation subgroup,
because \([X,(u,1)]=(du,1)\) and \(d\) is a unit.

Let \(A_f(U,V)\) be the Alexander coefficient of \(f\) relative
to the generator \([x,y]\) of the metabelian commutator module.
In that module the convention in the statement gives
\[
 [x,f^{-1}yf]
   =\bigl(1+(U-1)(V-1)A_f(U,V)\bigr)[x,y].
\]
The established diagonal compact-homology identification
therefore gives the full coefficient-two scalar
\begin{equation}
 b_r=1+d^2A_{f_r}(\zeta_{163},\zeta_{163}).
 \label{eq:163-alexander-scalar}
\end{equation}
This identity uses the representation on the full group \(G\).
On \(R\) it is linear over \(\mathcal O\), so
\(b_{r^n}=b_r^n\).

For a nonzero residue \(a\) modulo \(163\), write
\(d_a=\zeta_{163}^a-1\) and
\(\sigma_a(\zeta_{163})=\zeta_{163}^a\).
The power formula gives
\[
 X^a=(0,t^a),\qquad Y^a=(d^{-2}d_a,t^a),\qquad
 [X^a,Y^a]=(d^{-2}d_a^2,1).
\]
Consequently
\begin{equation}
 f(X^a,Y^a)
   =\bigl(d^{-2}d_a^2A_f(\zeta_{163}^a,\zeta_{163}^a),1\bigr).
 \label{eq:163-substitution}
\end{equation}
These evaluations all factor through the same finite quotient.
Explicitly,
\[
 \beta_a(v,t^i)=
       \bigl(d_a^2d^{-2}\sigma_a(v),t^{ai}\bigr)
\]
is an automorphism, \(\beta_a\beta_b=\beta_{ab}\), and
\(\beta_a(X)=X^a,\ \beta_a(Y)=Y^a\).

The scalar \(b_r\) need not be principal. Its norm is one,
and its odd Teichmuller factor has order dividing
\(2^{81}+1\), because \(c\) acts on that factor by the
\(2^{81}\)-power map. The chosen odd power thus gives
\begin{equation}
 b_{r^n}=b_r^n=h_{163}(r)^n
                     \equiv h_{163}(r)\pmod{4\mathcal O}.
 \label{eq:163-odd-power}
\end{equation}
It also preserves the image in \(E\), which has exponent two.
No identity between \(f_{r^n}\) and the word power \(f_r^n\)
is used.

Apply \eqref{eq:163-substitution} and
\eqref{eq:163-alexander-scalar} to \(f_{r^n}\).
The factors are translations, so their product has additive
coordinate
\[
 d^{-2}\sum_{j=0}^{26}s^{3j}(b_{r^n}-1)
   =2d^{-2}\sum_{j=0}^{26}s^{3j}
             \left((h_{163}(r)-1)/2\bmod2\mathcal O\right).
\]
This proves \eqref{eq:163-exact-word-value}. Multiplication
by \(2\) here denotes the canonical injection
\(\mathcal O/2\mathcal O\to2\mathcal O/4\mathcal O\);
there is no division by \(2\) in \(\mathcal O_4\).
That injection and multiplication by the unit \(d^{-2}\)
are injective. Equation \eqref{eq:163-trace-detector}
therefore proves \eqref{eq:163-word-equivalence}.
\end{proof}

The associated relative class has a unique normalized
representative:
\begin{equation}
 z(rh)=h_{163}(r)\bmod A,\qquad
 [z]\in\ker\bigl(H^1_{\rm cont}(G,E)\to
                         H^1_{\rm cont}(H,E)\bigr),\qquad z|_H=0.
 \label{eq:163-relative-twisted-class}
\end{equation}
Well-definedness follows from \(h_{163}(R\cap H)=A\);
equivariance gives the cocycle identity, and compact descent
from \(R\times H\) gives continuity.
If two cohomologous cocycles vanish on \(H\), their difference
is the coboundary of an element of \(E^H\). This fixed space
is zero, because the arithmetic cyclotomic \(C_3\)-action
is irreducible. Thus \([z]=0\) if and only if \(z=0\).
The common full image is \(A\) precisely when
\eqref{eq:163-word-equivalence} vanishes for every \(r\in R\).
Otherwise the quotient \(G\twoheadrightarrow E\rtimes C_3\simeq A_4\)
described above follows. This class uses \(G\) and the nontrivial
coefficient module \(E\); the character-extension torsor uses
\(G/R\) and \(E^\vee\).

The defining cycle relations alone have not yet been shown here
to force this universal word identity. Their diagonal
two-cycle specialization is automatic in \(\mathscr D\):
swapping \(X,Y\) acts as minus one on its commutator subgroup.
The three-cycle relation at multiplier one gives the necessary
Alexander identity
\[
 A_f(U,V)+V^{-1}A_f((UV)^{-1},U)
       +(UV)^{-1}A_f(V,(UV)^{-1})=0.
\]
Indeed \([z,x]=V^{-1}[x,y]\) and
\([y,z]=(UV)^{-1}[x,y]\) for \(z=(xy)^{-1}\).
Using two-cycle symmetry and then setting \(U=V=\zeta_{163}\)
gives
\[
 A_f(\zeta_{163},\zeta_{163})
   +(\zeta_{163}^{-1}+\zeta_{163}^{-2})
                 A_f(\zeta_{163},\zeta_{163}^{-2})=0.
\]
The second coefficient is a unit, so this equation relates
the diagonal and mixed values without making the displayed
trace vanish. The usual forgetful projection of the five-cycle
relation gives the two-cycle relation
\cite[Section~1, Proposition~1.3]{Furusho2003}.
A different use of the five-cycle relation may give more
information. No solution of a finite projected system is
being used as a realization of an element of \(G\).

Finite-field rigidity does not evaluate this detector. Guillot proves
that \(\mathrm{GT}_1(\mathrm{PSL}_2(2^s))\) is trivial, using
rigidity of generating triples with prescribed traces
\cite[Theorem~1.1, Proposition~3.2 and Corollary~3.3, pp.~1--3]{Guillot2016}.
Here the residual scalar lies in \(1+2O\), so its reduction modulo
two is already one. The test takes place modulo four. Applying
a finite-field rigidity theorem would therefore require an additional
integral lifting argument, which is not established here.
